\documentclass[UTF8]{ctexart}
\usepackage{amsmath,amssymb}
\usepackage[a4paper,margin=2.5cm]{geometry}
\title{数学分析 1014}
\date{}
\begin{document}
\maketitle

\section*{数学分析 1014}
\subsection*{0. Recall 夹迫原理}
$\exists \delta > 0 . \forall x . 0<|x-x_0|< \delta$ , 有$f(x) \leq g(x) \leq h(x)$, 若$\lim_{x \to x_0}f(x) = \lim_{x \to x_0}h(x) = A$, 则$\lim_{x \to x_0}g(x) = A$

e.g. 
\[
\lim\limits_{x \to x_0}g(x) = A
\]
回顾不等式
\[
\sin{x} \leq x \leq \tan{x}
\]
由$\sin{x} \leq x \ \Rightarrow \ \frac{\sin{x}}{x} \leq 1$
$x \leq \tan{x} = \frac{\sin{x}}{\cos{x}} \ \Rightarrow \ \cos{x} \leq \frac{\sin{x}}{x}$
我们就得到了$\cos{x} \leq \frac{\sin{x}}{x} \leq 1$我们希望$\cos{x}$趋于0的极限也是1,通过夹迫原理夹出极限的值
我们已知$\cos{x} = 1 - 2\sin{\frac{x}{2}}^2$
易得
\[
0 \leq |\cos{x}-1| \leq 2\sin{\frac{x}{2}}^2 \leq 2 \times (\frac{x}{2})^2 = \frac{1}{2}x^2
\]
夹迫定理易得
\[
\lim_{x \to 0} \cos{x} = 1
\]
即得
\[
\lim_{x \to 0^+} \frac{\sin{x}}{x} = 1
\]
接下来考虑$0^-$
注意到原函数为偶函数
或说明不等式在$(-\frac{\pi}{2} , 0)$也成立
可知x趋于0的左右极限均相等,均等于1

命题:
\[
\lim_{x \to x_0}f(x) = A \Leftrightarrow \lim_{x \to x_0^-}f(x) = \lim_{x \to x_0^+}f(x) = A
\]
证明不在此赘述
值得注意的是,此命题成立的根本原因在于$\mathbb{R}$是线性的,趋于零只有左右两个方式

e.g.
\[
\lim_{x \to +\infty}(1+\frac{1}{x})^x = e
\]
这次我们用数列极限来证明函数极限
不妨设x充分大 , 又已知$[x] \leq x \leq [x]+1$,由放缩易得
\[
(1+\frac{1}{[x]+1})^{[x]} \leq (1+\frac{1}{x})^x \leq (1+\frac{1}{[x]})^{[x]+1}
\]
此时两边已经可以套用数列极限,等价于
\[
\lim_{n \to +\infty}(1+\frac{1}{n+1})^n \qquad \lim_{n \to +\infty}(1+\frac{1}{n})^{n+1}
\]
由数列知识易得,两边都等于e,由夹迫定理可知$\lim_{x \to +\infty}(1+\frac{1}{x})^x = e$

\subsection*{1. 复合函数求极限}
假设$\lim_{u \to u_0}f(x) = A , \lim_{x \to x_0}g(x) = B$,$\exists x_0$的空心邻域, s.t. $g(x) \neq u_0$$\lim_{x \to x_0}f(g(x)) = A$

从外向里进行证明
证明:
\[
\lim\limits_{u \to u_0}f(x) = A \ \Leftrightarrow \ \forall \varepsilon > 0 , \exists u , 0<|u-u_0|<\eta , \text{有}|f(u)-A|<\varepsilon
\]
对于上面的$\eta >0$ , 由于$\lim_{x \to x_0}g(x) = u_0$

那么有
\[
\exists \delta > 0 , \ s.t. \ \forall x (0<|x-x_0|<\delta) , |g(x)-u_o|<\eta
\]
而又有g(x)在某个$x_0$的空心邻域里,$g(x) \neq u_0$
即得$0 < |g(x)-u_0| < \eta$
所以$f(g(x))-A|<\eta$
\begin{quote}
如果$\lim_{u \to u_0}f(x) = f(u_0)$,上面定理中高亮部分可以去掉
\end{quote}

\begin{quote}
$\lim_{u \to u_0}f(x) = f(u_0) \Leftrightarrow f(x)$在$u_0$处连续,而所有初等函数在定义域中都连续,故高亮部分一般不用验证
\end{quote}

\subsection*{2. 单调有界定理}
\noindent 定理:
\noindent $f$在$(x_0-\eta,x_0)$上单调上升且有上界,则$\lim_{x \to x_0^-}f(x)$存在\par
\noindent $f$在$(x_0,x_0+\eta)$上单调上升且有上界,则$\lim_{x \to x_0^+}f(x)$存在\par
\noindent 但是即使两个都满足,也不保证$\lim_{x \to x_0}$存在\par
\noindent 证明:
\noindent 取$\{ X_n \}$ s.t. $x_n<x_0$且$x_n \to x_0$严格递增,可知$\{ f(x_n) \}$单调递增且有上界
\noindent 可知$\lim_{n \to +\infty}$存在且为A
\noindent 下面证明$\lim_{x \to x_0^-}=A$
\[
\forall \varepsilon > 0 , \text{寻找} \delta > 0 , \forall x , x_0-\delta<x<x_0 , \text{有}A-\varepsilon<f(x) \leq A \\
\exists N , \forall n > N , \text{有}A-\varepsilon < f(x_n) \leq A
\]
\noindent 取$\delta = \frac{1}{2} (x_0-X_{N+1})$

\noindent 那么$\forall x , x_0-\delta<x<x_0 \Rightarrow x>x_{N+1}$
\noindent 则$f(x) > f(x_{N+1}>A- \varepsilon \Rightarrow f(x)>A- \varepsilon$
\noindent 又$x_n \to x_0$对于任意固定x,$\exists M , x_m>x$
\noindent 有$f(x) \leq f(x_M) \leq A$
\noindent 

\end{document}
